% ======================================================================
\section{Threat Model and Problem Statement}
\label{sec:problem}
% ======================================================================
We consider the network segment of a financial institution (for example
core banking, payment switching or online banking front-ends) monitored by a
flow exporter. Each bidirectional flow is summarized as a record
$\mathbf{r}=(\mathbf{r}^{\mathrm{num}},\mathbf{r}^{\mathrm{cat}})$ with $d$
numeric and $c$ categorical attributes. The adversary is an external or
insider attacker whose activity (reconnaissance, exploitation, DoS/DDoS,
brute force, botnet C\&C, web attacks, backdoors) produces flows. The task
is to output, for every flow, a threat score $p\in[0,1]$, a binary decision
$\hat{y}=\mathbb{1}[p\ge\tau]$, an attack family and a severity tier. A
second, \emph{honest-but-curious} party is also considered: an operator,
vendor or auditor who sees the detector's stored inputs and its released
preprocessing statistics. Against this party HALX must limit
re-identification of individual flows.

% ======================================================================
\section{Proposed Framework: HALX}
\label{sec:method}
% ======================================================================
Fig.~\ref{fig:architecture} shows the three stages: (1) a privacy-preserving
layer, (2) an Attention-LSTM encoder and (3) an XGBoost scoring head with
severity triage. Algorithm~\ref{alg:halx} summarizes training.

\begin{figure*}[t]
\centering
\resizebox{0.97\textwidth}{!}{% HALX architecture (TikZ). Included from main.tex.
\begin{tikzpicture}[
  font=\scriptsize, >=latex, node distance=3.2mm,
  blk/.style={draw, rounded corners=1.5pt, align=center, minimum height=6.5mm, inner sep=2pt, fill=#1},
  blk/.default=white,
  grp/.style={draw, dashed, rounded corners=3pt, inner sep=2.6pt},
  lab/.style={font=\tiny\itshape, text=black!70},
  arr/.style={->, thick, black!75}]
% ---- input and privacy layer
\node[blk=gray!12, minimum width=13mm] (raw) {Flow record\\$\mathbf{r}$ (raw)};
\node[blk=red!10, right=of raw, minimum width=17mm] (sup) {Suppress direct\\identifiers};
\node[blk=red!10, right=of sup, minimum width=19mm] (mond) {Mondrian\\$k$-anonymity on QIs};
\node[blk=red!10, right=of mond, minimum width=16mm] (dp) {$(\varepsilon,\delta)$-DP\\standardization};
\node[grp, draw=red!60, fit=(sup)(mond)(dp), label={[lab, anchor=south east]north east:(1) Privacy-preserving layer}] (P) {};
\draw[arr] (raw) -- (sup); \draw[arr] (sup) -- (mond); \draw[arr] (mond) -- (dp);
% ---- encoder
\node[blk=blue!8, below=8mm of raw.south west, anchor=north west, minimum width=13mm] (tok) {Feature\\tokenizer};
\node[blk=blue!8, right=of tok, minimum width=15mm] (lstm) {2-layer\\BiLSTM};
\node[blk=blue!8, right=of lstm, minimum width=15mm] (mha) {Multi-head\\self-attn.\,+\,LN};
\node[blk=blue!8, right=of mha, minimum width=13mm] (pool) {Attention\\pooling};
\node[grp, draw=blue!60, fit=(tok)(lstm)(mha)(pool), label={[lab, anchor=north west]south west:(2) Attention-LSTM encoder}] (E) {};
\draw[arr] (tok) -- node[above, lab]{$T{\times}d$} (lstm);
\draw[arr] (lstm) -- node[above, lab]{$\mathbf{H}$} (mha);
\draw[arr] (mha) -- node[above, lab]{$\mathbf{Z}$} (pool);
\draw[arr] (dp.south) |- ++(0,-3.2mm) -| node[pos=0.25, above, lab]{$\tilde{\mathbf{x}}$} (tok.north);
% ---- head
\node[blk=green!10, right=6mm of pool, minimum width=14mm] (cat) {$[\mathbf{e}\,\Vert\,\tilde{\mathbf{x}}]$};
\node[blk=green!10, right=of cat, minimum width=13mm] (xgb) {XGBoost\\$p(\text{attack})$};
\node[grp, draw=green!50!black, fit=(cat)(xgb), label={[lab, anchor=north east]south east:(3) Scoring head}] (X) {};
\draw[arr] (pool) -- node[above, lab]{$\mathbf{e}$} (cat);
\draw[arr] (cat) -- (xgb);
\draw[arr, dashed] (dp.east) -| (cat.north);
% ---- alerting
\node[blk=orange!15, right=6mm of xgb, minimum width=12mm] (alert) {Severity\\triage};
\draw[arr] (xgb) -- (alert);
\node[lab, below=0.5mm of alert, align=center] {low / med. / high\\+ top attended\\features};
\end{tikzpicture}
}
\caption{The HALX pipeline. Raw flow records pass through the privacy layer
(direct-identifier suppression, Mondrian $k$-anonymous generalization of
quasi-identifiers, DP standardization). The Attention-LSTM encoder reads the
privatized record as a sequence of feature tokens and produces an embedding
$\mathbf{e}$. XGBoost scores $[\mathbf{e}\Vert\tilde{\mathbf{x}}]$, and the score is
mapped to a severity tier together with the most-attended features.}
\label{fig:architecture}
\end{figure*}

\subsection{Privacy-Preserving Layer}
\label{sec:privacy}
\textit{Suppression.} Direct identifiers (IP addresses, ports,
timestamps, flow IDs) are dropped before any processing.

\textit{Quasi-identifier generalization.} Let $Q\subset\{1,\dots,d\}$ be the
quasi-identifiers (QIs). On UNSW-NB15 these are duration, source and
destination bytes, and source TTL; on CICIDS2017 they are flow duration,
forward and backward payload lengths, and the initial forward TCP window.
Volumes play the role of transaction sizes, and TTL or window values
fingerprint operating systems. After a public, data-independent transform
$g(x)=\operatorname{sign}(x)\log(1+|x|)$, we build a strict Mondrian
partition~\cite{lefevre2006mondrian} of the training QIs. Each region is
split recursively at the median of its widest normalized QI, as long as both
children keep at least $k$ records. Every record, in training or at
inference, is routed to its leaf and its QIs are replaced by the leaf median.
Hence every generalized QI vector the detector stores or logs is shared by
at least $k$ training records ($k$-anonymity~\cite{sweeney2002kanon}). We use
$k=10$ by default and verify the smallest equivalence class after fitting.

\textit{DP standardization.} Standardization needs per-feature means and
variances, which are normally shipped with a model and can reveal individual
records when an outlier dominates. We release them with the Gaussian
mechanism. With values clipped to $[-C,C]$ ($C=25$ in $g$-space) and public
$n$, the sums $S_1=\sum_i x_i$ and $S_2=\sum_i x_i^2$ have replace-one
sensitivities $2C$ and $C^2$. Each of the $2d$ released sums receives noise
$\mathcal{N}(0,\Delta^2/2\rho_q)$ with $\rho_q=\rho/2d$, where
\begin{equation}
  \rho=\Bigl(\sqrt{\ln(1/\delta)+\varepsilon}-\sqrt{\ln(1/\delta)}\Bigr)^2
  \label{eq:rho}
\end{equation}
is the largest $\rho$ for which $\rho$-zCDP implies
$(\varepsilon,\delta)$-DP, using $\varepsilon=\rho+2\sqrt{\rho\ln(1/\delta)}$~\cite{bun2016zcdp}.
By composition, the released statistics satisfy
$(\varepsilon,\delta)$-DP. We use $\varepsilon=1$ and $\delta=10^{-5}$.
\emph{Scope:} this protects the shipped statistics only. The network
weights are not differentially private; that would require
DP-SGD~\cite{abadi2016dpsgd} (Sec.~\ref{sec:discussion}).

\subsection{Attention-LSTM Encoder}
\label{sec:encoder}
UNSW-NB15's partition files carry no timestamps and are sorted by label, and
the CICIDS2017 subset is randomly sampled. Sliding windows of
\emph{consecutive flows} would therefore leak labels through ordering, a form of
data snooping~\cite{arp2022dos}. HALX instead reads each flow record as a
sequence of $T=d+c$ \emph{feature tokens}, in the spirit of tabular
transformers~\cite{gorishniy2021revisiting}:
\begin{equation}
  \mathbf{t}_j=\tilde{x}_j\,\mathbf{w}_j+\mathbf{b}_j\ (j\le d),\qquad
  \mathbf{t}_{d+i}=\mathbf{E}_i[r^{\mathrm{cat}}_i]+\mathbf{b}_{d+i},
  \label{eq:tok}
\end{equation}
with $\mathbf{w}_j,\mathbf{b}_j\in\mathbb{R}^{16}$. A two-layer bidirectional
LSTM~\cite{hochreiter1997lstm} with 32 units per direction produces
$\mathbf{H}=[\mathbf{h}_1,\dots,\mathbf{h}_T]\in\mathbb{R}^{T\times64}$. Here
the recurrence models interactions between related features, for example
packet counts, sizes and rates of the same direction. Multi-head
self-attention~\cite{vaswani2017attention} with $H=4$ heads lets every token
attend to every other:
\begin{equation}
  \mathbf{Z}=\mathrm{LN}\bigl(\mathbf{H}+\mathrm{MHA}(\mathbf{H}\mathbf{W}^Q,
  \mathbf{H}\mathbf{W}^K,\mathbf{H}\mathbf{W}^V)\bigr),
  \label{eq:mha}
\end{equation}
where each head computes
$\mathrm{softmax}(QK^{\top}/\sqrt{d_k})V$. Additive attention pooling
then yields the flow embedding together with an importance distribution
over features:
\begin{equation}
  \alpha_j=\operatorname{softmax}_j\bigl(\mathbf{u}^{\top}\tanh(\mathbf{W}_a\mathbf{z}_j)\bigr),\qquad
  \mathbf{e}=\textstyle\sum_{j}\alpha_j\mathbf{z}_j .
  \label{eq:pool}
\end{equation}
The weights $\alpha_j$ are reported with each alert as an explanation.
The encoder is trained end to end with a linear softmax head on the
multi-class labels (cross-entropy, AdamW, learning rate $3\cdot10^{-3}$,
cosine schedule, batch 1024). Early stopping uses validation macro-F1.

\subsection{XGBoost Scoring Head and Severity Triage}
\label{sec:head}
The deep encoder captures cross-feature structure, but on tabular data
boosted trees model sharp, axis-aligned decision boundaries better~\cite{grinsztajn2022tree}.
The head therefore receives the concatenation
$\mathbf{f}=[\mathbf{e}\,\Vert\,\tilde{\mathbf{x}}]$ of the embedding and the
privatized features. One XGBoost model~\cite{chen2016xgboost} estimates
$p=P(\text{attack}\mid\mathbf{f})$ and a second predicts the attack family
(up to 800 trees, depth 8, learning rate 0.08, early stopping on the
validation split). Alerts with $p\ge\tau=0.5$ are triaged into
\emph{low} ($p<0.8$), \emph{medium} ($0.8\le p<0.95$) and \emph{high}
($p\ge0.95$) severity, so SOC analysts can work through high-severity alerts first.

\begin{algorithm}[t]
\caption{HALX training}
\label{alg:halx}
\begin{algorithmic}[1]
\REQUIRE training/validation flows $\mathcal{D}_{tr},\mathcal{D}_{va}$; QI set $Q$; $k,\varepsilon,\delta$
\STATE Drop direct identifiers; apply $g(\cdot)$; impute with training medians
\STATE Fit Mondrian tree on $\mathcal{D}_{tr}[Q]$ with leaf size $\ge k$; generalize $Q$
\STATE Release $(\mu,\sigma)$ via Gaussian mechanism with $\rho(\varepsilon,\delta)$ (Eq.~\ref{eq:rho}); standardize $\to\tilde{\mathbf{x}}$
\STATE Oversample minority families in $\mathcal{D}_{tr}$ with SMOTE-NC~\cite{chawla2002smote}
\STATE Train encoder (Eqs.~\ref{eq:tok}--\ref{eq:pool}) with early stopping on $\mathcal{D}_{va}$
\STATE Extract $\mathbf{f}=[\mathbf{e}\Vert\tilde{\mathbf{x}}]$ for $\mathcal{D}_{tr},\mathcal{D}_{va}$
\STATE Fit binary and multi-class XGBoost on $\mathcal{D}_{tr}$, early-stopped on $\mathcal{D}_{va}$
\RETURN privacy layer, encoder, XGBoost heads, threshold $\tau$
\end{algorithmic}
\end{algorithm}
